\documentclass[twocolumn, 10pt]{article}
\usepackage[utf8]{inputenc}
\usepackage[spanish]{babel}
\usepackage{amsmath, amssymb, amsthm}
\usepackage{algorithm}
\usepackage{algpseudocode}
\usepackage{listings}
\usepackage{xcolor}
\usepackage{geometry}
\usepackage{hyperref}
\usepackage{booktabs}
\usepackage{graphicx}

\geometry{
  top=1.8cm, bottom=1.8cm,
  left=1.5cm, right=1.5cm,
  columnsep=0.6cm
}

\definecolor{codebg}{rgb}{0.95,0.95,0.95}
\definecolor{codekey}{rgb}{0.0,0.35,0.65}
\definecolor{codecomment}{rgb}{0.3,0.5,0.3}

\lstset{
  backgroundcolor=\color{codebg},
  basicstyle=\ttfamily\fontsize{7}{8.5}\selectfont,
  keywordstyle=\color{codekey}\bfseries,
  commentstyle=\color{codecomment}\itshape,
  breaklines=true,
  frame=single,
  rulecolor=\color{gray!40},
  xleftmargin=4pt, xrightmargin=4pt,
  aboveskip=4pt, belowskip=4pt
}

\newtheorem{theorem}{Teorema}

\title{\textbf{Exponenciación Modular Rápida:\\
De $O(n)$ a $O(\log n)$ y su Rol en la Criptografía RSA}\\[4pt]
\small UNIVALLE BUGA -- Análisis y Diseño de Algoritmos -- 2026}

\author{Jhon Alejandro Ortega Valdes\\ Julian Andres Martinez Torres}
\date{}

\begin{document}
\maketitle

%% ─────────────────────────────────────────
\section*{Resumen}
%% ─────────────────────────────────────────

Calcular $a^n \bmod m$ para valores enormes de $n$ es la operación central
de la criptografía RSA, donde los exponentes superan los 2048 dígitos
binarios. El método directo ---multiplicar $n$ veces--- tiene un costo
$O(n)$, es decir, el tiempo crece al mismo ritmo que el exponente, lo cual
lo vuelve completamente inviable para esas magnitudes. La técnica de
\emph{exponenciación por cuadrados sucesivos}, formalizada por D.\ Knuth
(1969) y basada en ideas de Gauss (1801), reduce ese costo a $O(\log n)$:
el tiempo crece muchísimo más despacio que el exponente, convirtiendo una
operación imposible en algo que tarda microsegundos. Este documento
demuestra esa reducción, muestra implementaciones en Python, y la
valida con mediciones reales de tiempo.

%% ─────────────────────────────────────────
\section{Introducción}
%% ─────────────────────────────────────────

\textbf{¿Qué significa calcular $a^n \bmod m$?} El operador
\textbf{módulo} ($\bmod$) entrega el \emph{resto} de una división. Por
ejemplo:
\[
  17 \bmod 5 = 2 \quad \text{(porque } 17 = 3 \times 5 + 2\text{)}
\]

La \textbf{exponenciación modular} es entonces elevar $a$ a la potencia
$n$ y quedarse con el resto al dividir entre $m$:
\[
  r = a^n \bmod m
\]

\textbf{Ejemplo.} Un reloj de 12 horas aplica el módulo de forma
continua sin que el usuario lo note. Si en un momento dado el reloj
marca las 10~h y transcurren 5 horas más, el resultado no es 15~h
sino $15 \bmod 12 = 3$~h: el reloj ``reinicia'' al llegar a 12 y
continúa desde cero. Lo mismo ocurre si se suman 29 horas a las 10~h:
$39 \bmod 12 = 3$~h, exactamente el mismo resultado que sumar solo 3.
Esta propiedad ---que valores muy distintos pueden producir el mismo
resto--- es precisamente lo que hace útil el módulo en criptografía,
pues permite ocultar información dentro de operaciones que son fáciles
de calcular en una dirección pero muy difíciles de revertir. RSA
aprovecha ese mismo principio, pero con módulos del orden de
$10^{600}$ en lugar de 12.

En el cifrado RSA, proteger un mensaje $M$ se hace calculando
$c = M^e \bmod n$, y recuperarlo requiere $M = c^d \bmod n$, donde $d$
puede ser del orden de $2^{2048}$. Ejecutar $2^{2048}$
multiplicaciones convencionales tomaría aproximadamente $10^{597}$ años
incluso con los computadores más rápidos del mundo. La exponenciación por
cuadrados sucesivos resuelve esto con apenas $2048$ multiplicaciones.

%% ─────────────────────────────────────────
\section{Análisis de Complejidad}
%% ─────────────────────────────────────────

\subsection{Método directo: \texorpdfstring{$O(n)$}{O(n)}}

La idea más simple es multiplicar $a$ por sí mismo, una y otra vez,
exactamente $n$ veces:

\[
  a^n = \underbrace{a \cdot a \cdot \ldots \cdot a}_{n \text{ veces}}
\]

\textbf{Ejemplo.} Un repartidor que debe entregar un paquete en el
décimo piso de un edificio sin ascensor sube los escalones de uno en
uno: cada escalón es un paso obligatorio y no puede saltarse ninguno.
Si el edificio tuviera 1000 pisos, daría 1000 pasos; si tuviera
$10^{616}$ pisos, el tiempo de subida sería incompatible con cualquier
expectativa de vida. El método directo de exponenciación funciona igual:
por cada unidad que crece el exponente $n$, el programa debe ejecutar
una multiplicación adicional. El costo es estrictamente proporcional
al tamaño del problema.

\begin{algorithm}[H]
\caption{Exponenciación directa (lenta)}
\begin{algorithmic}[1]
\Function{PotenciaDirect}{$a, n, m$}
  \State $r \gets 1$
  \For{$i = 1$ \textbf{hasta} $n$}
    \State $r \gets (r \cdot a) \bmod m$  \Comment{una multiplicación por paso}
  \EndFor
  \State \Return $r$
\EndFunction
\end{algorithmic}
\end{algorithm}

\textbf{Costo:} $n$ multiplicaciones $\Rightarrow T(n) = O(n)$. El tiempo
crece en proporción directa al exponente.

\subsection{Método rápido: \texorpdfstring{$O(\log n)$}{O(log n)}}

La clave está en una observación sencilla: en lugar de multiplicar
paso a paso, podemos \emph{doblar} el resultado en cada ronda.

\textbf{Ejemplo.} Una hoja de papel tiene un grosor de aproximadamente
0{,}1~mm. Al doblarla una vez, el grosor se duplica: 0{,}2~mm. Al
doblarla dos veces, 0{,}4~mm. Tras 10 dobleces se obtiene un grosor de
$2^{10} \times 0{,}1 = 102{,}4$~mm, es decir, más de 10~cm. Con el
método directo, para alcanzar esas 1024 capas habría que apilar 1024
hojas una por una: 1024 operaciones. Con el método de los dobleces,
solo se necesitan 10 operaciones para llegar al mismo resultado. Ese
es exactamente el principio de la exponenciación por cuadrados
sucesivos: en lugar de acumular de a uno, se dobla el resultado en
cada paso, reduciendo drásticamente la cantidad de operaciones
necesarias.

La idea se basa en la siguiente propiedad:
\[
  a^n =
  \begin{cases}
    1 & \text{si } n = 0 \\
    \left(a^{n/2}\right)^2 & \text{si } n \text{ es par} \\
    a \cdot a^{n-1} & \text{si } n \text{ es impar}
  \end{cases}
\]

\textbf{Ejemplo concreto.} Para calcular $3^{13}$:
\begin{itemize}
  \item $13$ es impar: $3^{13} = 3 \cdot 3^{12}$
  \item $12$ es par: $3^{12} = (3^6)^2$
  \item $6$ es par: $3^{6} = (3^3)^2$
  \item $3$ es impar: $3^{3} = 3 \cdot 3^{2}$
  \item $3^2 = 9$, luego $3^3=27$, $3^6=729$, $3^{12}=531441$, $3^{13}=1594323$
\end{itemize}
Solo se necesitaron \textbf{5 multiplicaciones} en lugar de 12.

\begin{algorithm}[H]
\caption{Exponenciación por cuadrados sucesivos (rápida)}
\begin{algorithmic}[1]
\Function{PotenciaRapida}{$a, n, m$}
  \State $r \gets 1$
  \State $a \gets a \bmod m$
  \While{$n > 0$}
    \If{$n$ es impar}
      \State $r \gets (r \cdot a) \bmod m$
    \EndIf
    \State $n \gets \lfloor n/2 \rfloor$  \Comment{reducir a la mitad}
    \State $a \gets (a \cdot a) \bmod m$  \Comment{cuadrar la base}
  \EndWhile
  \State \Return $r$
\EndFunction
\end{algorithmic}
\end{algorithm}

\subsubsection{Demostración del costo}

\begin{theorem}
\textsc{PotenciaRapida} realiza a lo sumo $2\lfloor \log_2 n \rfloor + 2$
multiplicaciones, por lo que su costo es $T(n) = O(\log n)$.
\end{theorem}

\begin{proof}
El ciclo \textbf{mientras} repite su cuerpo mientras $n > 0$. En cada
vuelta, $n$ se divide entre 2. La cantidad de veces que se puede dividir
$n$ entre 2 antes de llegar a 0 es exactamente:
\[
  k = \lfloor \log_2 n \rfloor + 1
\]
En cada vuelta se hace \emph{siempre} 1 cuadrado, y \emph{a lo sumo}
1 multiplicación extra cuando $n$ es impar. Por eso:
\[
  \text{total de multiplicaciones} \leq 2k = 2(\lfloor \log_2 n \rfloor + 1)
\]
Como $\lfloor \log_2 n \rfloor < \log_2 n$, el costo es $T(n) = O(\log n)$.
\qed
\end{proof}

\subsection{Comparación de operaciones}

\begin{table}[H]
\centering
\caption{Multiplicaciones necesarias según el método}
\small
\begin{tabular}{@{}lcc@{}}
\toprule
Exponente $n$ & Directo $O(n)$ & Rápido $O(\log n)$ \\
\midrule
$10^3$     & 1.000             & 10 \\
$10^6$     & 1.000.000         & 20 \\
$10^{12}$  & $10^{12}$         & 40 \\
$2^{2048}$ & $\approx 10^{616}$ & 2.048 \\
\bottomrule
\end{tabular}
\end{table}

La diferencia ya no es solo de velocidad: es la diferencia entre
lo que es posible y lo que no lo es.

%% ─────────────────────────────────────────
\section{Aplicaciones Reales}
%% ─────────────────────────────────────────

\subsection{Cifrado RSA}

RSA (Rivest, Shamir, Adleman, 1977) es el sistema de cifrado asimétrico
más utilizado en el mundo. Se basa en que multiplicar dos números primos
grandes es fácil, pero factorizarlos (encontrar de cuál multiplicación
surgieron) es extremadamente difícil. Sus pasos son:

\begin{enumerate}
  \item Elegir dos números primos grandes $p$ y $q$; calcular $n = p \cdot q$.
  \item Calcular $\varphi(n) = (p-1)(q-1)$, que cuenta cuántos números
        menores que $n$ no comparten factores con él.
  \item Elegir $e$ tal que $\text{mcd}(e,\, \varphi(n)) = 1$
        (el máximo común divisor sea 1, es decir, sean coprimos).
  \item Hallar $d = e^{-1} \bmod \varphi(n)$ (el inverso modular de $e$).
  \item \textbf{Cifrar:} $c = M^e \bmod n$ \quad (necesita \textsc{PotenciaRapida})
  \item \textbf{Descifrar:} $M = c^d \bmod n$ \quad (necesita \textsc{PotenciaRapida})
\end{enumerate}

\textbf{Ejemplo simplificado.} Con $p=61$, $q=53$, $n=3233$, $e=17$ y
$d=2753$: el mensaje $M=42$ se cifra como $42^{17} \bmod 3233 = 2557$, y
se descifra como $2557^{2753} \bmod 3233 = 42$. Ambas operaciones usan
\textsc{PotenciaRapida} y tardan microsegundos.

Sin $O(\log n)$, protocolos como HTTPS, las contraseñas de correo
electrónico y las transacciones bancarias en línea serían inviables.

\subsection{Otras aplicaciones}

\textbf{Acuerdo de claves Diffie-Hellman (1976):} Permite que dos personas
acuerden una clave secreta a través de un canal público, como si hablaran
en clave delante de todos. Calcula $g^a \bmod p$ con valores de $a$
mayores a $10^{300}$, por lo que también depende de la exponenciación
rápida.

\textbf{Prueba de primalidad de Miller-Rabin:} Determina si un número es
primo de forma eficiente. Se usa para generar las claves de RSA y requiere
la exponenciación rápida en su núcleo.

\textbf{Criptografía de curva elíptica:} Variante moderna que aplica el
mismo principio de cuadrados sucesivos sobre puntos geométricos en lugar
de números enteros. Logra la misma seguridad que RSA-3072 usando claves
de solo 256 dígitos binarios, por lo que es el estándar en teléfonos
móviles y dispositivos con recursos limitados.

%% ─────────────────────────────────────────
\section{Resultados de las Pruebas}
%% ─────────────────────────────────────────

Se midió el tiempo de ejecución en Python 3 promediando 100 repeticiones
por valor de $n$, sobre un procesador Intel Core i5 a 3,5~GHz.
La aceleración indica cuántas veces más rápido es el
método rápido frente al directo.

\begin{table}[H]
\centering
\caption{Tiempo medido en Python 3 (ms), promedio de 100 repeticiones}
\small
\begin{tabular}{@{}lrrr@{}}
\toprule
$n$ & Directo (ms) & Rápido (ms) & Aceleración \\
\midrule
$10^3$             & 0,118  & 0,00262 & $45\times$ \\
$5\times10^3$      & 0,400  & 0,00198 & $202\times$ \\
$10^4$             & 0,797  & 0,00210 & $379\times$ \\
$5\times10^4$      & 4,131  & 0,00239 & $1.732\times$ \\
$10^5$             & 8,030  & 0,00293 & $2.742\times$ \\
$10^6$             & inviable & 0,00350 & --- \\
\bottomrule
\end{tabular}
\end{table}

\noindent Las gráficas completas y los archivos de prueba se encuentran
en el archivo \texttt{comparador\_rsa.py} del repositorio.

%% ─────────────────────────────────────────
\section{Conclusión}
%% ─────────────────────────────────────────

La exponenciación por cuadrados sucesivos es un ejemplo perfecto de cómo
un ingenio matemático sencillo ---en lugar de multiplicar $n$ veces,
dividir el problema a la mitad en cada paso--- transforma lo imposible en
trivial. Pasar de $O(n)$ a $O(\log n)$ significa que para $n = 2^{2048}$,
el número de operaciones baja de $10^{616}$ a apenas 2048: la misma tarea,
un abismo de diferencia. Cada vez que alguien envía un mensaje cifrado,
realiza una compra en línea o se conecta a una red segura, este algoritmo
---cuyas raíces se remontan al trabajo de Gauss en 1801--- opera en
segundo plano haciendo posible la vida digital moderna.

%% ─────────────────────────────────────────
\section*{Repositorio}
%% ─────────────────────────────────────────

\noindent
Código fuente y pruebas de rendimiento:\\
\url{https://github.com/Alejandro-Ortega10/Algoritmo-RSA.git}

\end{document}
